\section{Texto e organização}

\subsection{Ênfase, destaque e nota de rodapé}

\begin{lstlisting}[language={[LaTeX]TeX},caption={Formatação de texto}]
\textbf{negrito}
\textit{italico}
\emph{enfase conforme o contexto}
\underline{sublinhado}
\hl{texto destacado}
Texto com uma nota\footnote{Conteudo da nota.}
\end{lstlisting}

O resultado é: \textbf{negrito}, \textit{itálico}, \emph{ênfase conforme o
contexto}, \underline{sublinhado} e \hl{texto destacado}. Também é possível
acrescentar uma nota de rodapé ao fim de uma frase.\footnote{Este é o conteúdo da
nota de rodapé criada com \texttt{\textbackslash footnote}.}

\subsection{Listas}

Use \texttt{itemize} para uma lista sem ordem e \texttt{enumerate} quando a ordem
dos itens for importante.

\begin{lstlisting}[language={[LaTeX]TeX},caption={Listas com e sem numeração}]
\begin{itemize}
  \item Item sem numeracao.
  \item Outro item.
\end{itemize}

\begin{enumerate}
  \item Primeiro item.
  \item Segundo item.
\end{enumerate}
\end{lstlisting}

\begin{itemize}
  \item Item sem numeração.
  \item Outro item.
\end{itemize}

\begin{enumerate}
  \item Primeiro item.
  \item Segundo item.
\end{enumerate}

\subsection{Citação longa}

O ambiente personalizado \texttt{longcite} aplica fonte menor, espaçamento
simples e recuo ao trecho.

\begin{lstlisting}[language={[LaTeX]TeX},caption={Citação longa}]
\begin{longcite}
Trecho de uma citacao com mais de tres linhas. Ao final, informe autor,
ano e pagina de acordo com a norma adotada.
\end{longcite}
\end{lstlisting}

\begin{longcite}
Este é um exemplo visual de citação longa. Em um relatório real, substitua este
texto pelo trecho consultado e informe autor, ano e página. Evite copiar conteúdo
sem a indicação adequada da fonte.
\end{longcite}
